% 03-core.tex: the credential core.
\section{The core: one credential, one signature, one person}
\label{sec:core}

\motif{Privacy is not freedom from coercion.}

\subsection{What the credential is}

At the centre of Polaris is one row of the table \entity{IdentityToken} and at least one row of
\entity{TokenSignature}. The token row records a unique \code{token\_value}, a physical serial, the
person it belongs to, the agency that issued it, the algorithm it was issued under, its predecessor if
it replaced an earlier token, and its status. The signature row records the issuing agency's ML-DSA
signature over the SHA3-256 digest of the token value, together with the public key that produced
it, so that verification is self-contained and survives the retirement of the signing key. A token
without a signature row cannot exist: a database trigger refuses any write that would leave one
without a non-deprecated signature. The credential is therefore not a document but a signed claim:
\emph{this authority issued this token to this person.} The claim may also bind a
key the holder controls, recorded in \entity{HolderKeyEvent} under the same append-only discipline
as an authority key, so that possession of the file stops being possession of the credential
(Section~\ref{sec:agents}). \sImpl

What the holder holds is the authenticity pack: the token value, the signature and the issuer's
public key, exported by the application and kept as a file by the reference wallet
(\code{scripts/polaris-wallet.py}). Anyone can check the pack with a standard ML-DSA library and no
Polaris code (Section~\ref{sec:verifiers}).

\subsection{The physical token, as an object somebody else can implement}

The silicon is modelled and not manufactured; everything around it exists. The card is an encoding with a
reference codec and published vectors (\code{polaris\_card/\allowbreak{}card\_profile.py}), a software token that
speaks ISO 7816-4 with real status words, two PIN slots, a retry counter and a PUK
(\code{emulator.py}), a personalization step in which the card generates its own keys so that the
authority never holds them, and a reference verifier device that reads a card over NFC or a
presentation over QR and reports three findings separately, possession, authenticity and
authorization, because a revoked card still signs a challenge and a device that says only
\emph{accepted} teaches its operator nothing. The profile states its one shaping constraint plainly:
certified secure elements that sign FIPS 204 in hardware are not something an authority can buy in
volume today, so a card carries two issuer signatures over one body, a classical one that today's
silicon can produce and an ML-DSA one, and the reader decides which it trusts. Every one of these
runs in a drill on every push. \sImpl{} for the profile, the emulator, the vectors and the device;
\sFut{} for a manufactured card, which no drill can establish.

\figfile{nucleus}{The core entity geometry, unchanged from Version 1. Seven entities around
\entity{IdentityToken}: the person who holds it, the agency that issues it, the algorithm that signs
it, the contexts it may be verified in, and the two append-only event logs it generates. The
many-to-many relationships resolve to the junction tables \entity{AgencyAlgorithmAuth} and
\entity{TokenPermission}. This geometry is the constant of the paper: Section~\ref{sec:walls} draws
the other thirty-eight tables as rings around it.}{fig:nucleus}

\subsection{What the credential is not}

The constitution's tenth constraint, C10, states that identity attestation never carries spending
authority: there is no monetary claim in the schema, and a check fails the build if one appears. The
reason is stated as a failure mode, not a preference. A credential that is also money inherits
programmability: every constraint anyone wants to place on spending accretes onto the identity token,
until an administrative paperwork error becomes an existential bank-balance error, and until the
system can be told that a person may not buy fuel. The same refusal covers reputation, behaviour and
profiles. Polaris answers one question, \emph{is this token valid for this context right now}, and
it declines to answer \emph{should this person be allowed to do X}. That second question belongs to
the relying party, outside the boundary. \sImpl

\subsection{One active token per person}

A person may hold several provisioned tokens but only one may be \code{ACTIVE}. Two active tokens for
one person would let one token authorise a transaction and the other repudiate it. The rule is C3,
and it is enforced by a partial unique index, \code{uq\_one\_active\_per\_person}, on the person's
identifier where the status is active. Two operators activating two reserve tokens for the same
holder at the same moment find that exactly one of them succeeds, and the test that proves it runs
with real threads against a live database (C9). The rule is also modelled in TLA+ and the model is
checked on every push, under every interleaving of concurrent issue and revoke, with a counterpart
configuration it is required to fail. One more thing is true of it, recorded in \evidencerecord: when the index was dropped to see what would notice, the suite that exists to prove
C3 turned out to be violating it, committing a second active token before failing. \sImpl

What C3 calls a person is a row, \code{Individual.individual\_id}, not a human being. The index makes
a second active credential for the same row impossible; it does nothing about the same human
enrolled as two rows, and no mechanism in the repository does. Uniqueness of humans rests on
enrolment, where an operator records the evidence and the assurance level is derived from it.

\subsection{The lifecycle}

\figfile{lifecycle}{The token lifecycle machine as the trigger \code{trg\_token\_state\_machine}
enforces it: six legal transitions, each labelled with the lifecycle event it writes automatically.
\code{ACTIVE} is the only operational state. A token enters at \code{RESERVE} and is activated, or
is revoked before it ever is; an active token leaves through one of three terminal states or into
\code{DORMANT} when a successor takes over.}{fig:lifecycle}

A trigger enforces the transitions, and a second trigger writes a \entity{TokenLifecycleEvent} row for every status
change, including a change made by hand from a database shell. Succession is by reference, never by
overwrite: a replacement token points at its predecessor and the old token stays in the database in
its terminal state. Lost tokens stay lost, and their data is not erased. \sImpl

\subsection{Entering, leaving and returning}

Three mechanisms surround the core so that the schema cannot be turned against the holder at the
edges of the lifecycle, and each keeps a record of what it rested on.

\begin{description}
  \item[Entering: tiered enrolment, and what an enrolment rests on.] Every person has an enrolment
    status drawn from five values, and \code{EXEMPT} is a first-class one: a recognised way to take
    part in civic life without a token, whether for biometric incompatibility, religious exemption or
    an alternative path. Recording an exemption is one row; enumerating the not-enrolled is possible
    but unhelped by any function, and it is audited when done. A credential records how the person
    was proven to be who they claimed: \entity{EnrollmentProofing} records one proofing event and \entity{EnrollmentEvidence} what
    it rested on, both append-only, and the assurance level is \emph{derived} from the evidence
    rather than typed by an operator; there is deliberately no column for the document itself, and a
    named set of fields is refused outright, so the record cannot reconstruct a passport. What it
    does not establish is stated in the review packet as a limitation: whether the evidence was
    genuine is taken on the operator's word. \sImpl

    Every combination of evidence a standard names starts from documents, and a person with none, a
    care leaver, a refugee, an adult who never held a passport, fails all of them by omission. The
    trusted referee (\entity{RefereeVouching}) is the answer to that, and it is also the easiest way
    to mint an assurance level out of nothing, so the repository treats the anti-exclusion mechanism
    and the forgery channel as the same table: a vouching is bounded, co-signed past a threshold, and
    invisible on the credential itself, because a person who needed one should not carry a mark for
    it at every counter. An enrolment code, the last piece that needs no hardware, establishes only
    that somebody who could reach a channel returned a secret, and the level it can contribute is
    capped accordingly. \sImpl
  \item[Leaving: issuer discretion, as a decision that survives being changed.] The worst failure of
    an identity system is not a forged token but an authority revoking a population's tokens at
    scale. A trigger bounds how fast one agency can revoke its own tokens (five percent of its
    outstanding base in thirty days by default) and above that a co-signer from a different agency is
    required, recorded in the audit trail. The five percent is a deployment default and a governance
    parameter, not a security invariant: an authority may set its own (one seeded agency runs at
    seven), and no property of the system follows from the number itself. The bound converts a surprise into a trend that reporting
    can see; it does not stop slow abuse under the ceiling, and says so. The settings around this
    power are records: the ceiling itself (\entity{IssuerDiscretionPolicy}) can
    only be changed with a stated justification and leaves a row saying who changed it and from what;
    a quota on how much of the population an agency may touch is a decision with a reason rather than
    a value that can be overwritten; the row at the root of the hierarchy, \entity{Agency}, carries a
    trigger that records every change to an authority's level and marks whether it widened; and a
    relying party's bar is recorded with the same discipline, weakenings flagged and reasons required. \sImpl
  \item[Returning: the recovery ceremony.] A fire takes the wallet and the reserve together. Recovery
    is mandatory and is also an impersonation path, so it is made expensive without being made
    impossible: a two-phase ceremony with a 48-hour cool-down, three independent out-of-band channels,
    an approver who cannot be the requester, and admin-only completion. The decision fields move one
    way under a trigger, and the four-eyes rule, the cool-down and the three channels are refusals a
    stored procedure makes, each of which a mutation drill shows the tests would notice losing
    (Section~\ref{sec:walls}). No path in the product records the three channels
    (Section~\ref{sec:limitations}). The holder is without identity for the cool-down; the
    substitute credential a production system would want for that window is not built. \sImpl{}
    for the ceremony, \sFut{} for the interim credential.
\end{description}

\subsection{Duress: an alert the coercer cannot see}

The vocation above the ten constraints is anti-coercion: no person may be compelled to renounce,
transfer or surrender their identity against their will. The mechanism that gives that sentence
something to stand on is old and comes from banking: a second secret that looks exactly like
success to whoever is watching, while an alert goes out through a channel the watcher cannot see.
A holder may enrol a duress code. When it is presented, the verification proceeds identically on
every operator-visible surface, the same page, the same flash message, the same event row, and a
\entity{DuressEvent} row is appended to a table the operator's screens never join, from which a
pager fires at the highest severity. The comparison is constant-time and the work is paid on both
paths. On the card the same mechanism is a second PIN of the same length as the first, entered the
same way, unlocking a second key slot whose public key the card object does not carry, and a drill
measures the response in time as well as in bytes, because a coercer at a reader observes how long
the card took. \sImpl{} as a tested property of the modelled flow.

The duress mechanism is \emph{duress-aware}, not compulsion-resistant: it can keep an alert from a
coercer who does not know it exists, and it does not protect against an informed coercer or against
institutional access. A check refuses the stronger word on every outward surface. Against a casual coercer, one who does
not know duress codes exist, the mechanism works: the screen is identical and the alert is silent.
Against an informed coercer standing beside the holder it does not: the code is typed on the
verifier's terminal, keystrokes are not constant-time, and because enrolment is opt-in, \emph{I have no
duress code} is a claim a coercer can press on and the holder cannot prove. Against post-hoc lawful or
institutional access it is worse than nothing: \entity{DuressEvent} is append-only with no purge path,
which is C1 working as intended, so where the coercing party is the institution, or can later compel
it, the mechanism manufactures indelible evidence that the holder resisted. The assessment proposes
no mechanism change and this paper proposes none; it records the boundary. The side-channel claim,
that nothing observable distinguishes the two paths, has been measured by the repository and by
nobody else. \sExt

\begin{layercard}{Layer card: the credential core}
\qa{What it is}{One \entity{IdentityToken} row and its \entity{TokenSignature} rows, optionally bound to a holder key: a signed claim that an authority issued this token to this person, in a lifecycle with one active state, with the proofing it rested on recorded beside it.}
\qa{Why it exists}{Every other structure in Polaris organises itself around this claim. It is kept as small as possible so that the rest can be bounded.}
\qa{What it trusts}{The issuing agency's signing key, held behind a custody interface (Section~\ref{sec:signatures}); the schema's constraints; the operator's word that presented evidence was genuine.}
\qa{What it refuses}{Application code as a source of truth: the state machine, uniqueness and audit are decided by the database. An assurance level typed rather than derived. Money, reputation and profiles: C10.}
\qa{What it can see}{The holder's legal name, date of birth and country; the issuing agency; the algorithm; the lifecycle; the kind of evidence an enrolment rested on.}
\qa{What it cannot see}{A biometric template (never stored), the document behind an enrolment (no column exists for it), a verification graph (Section~\ref{sec:privacy}), a spending history (never modelled).}
\qa{If it fails}{A forged core is a forged signature, which the two-witness verification and the detached verifier refuse; a duplicated core is refused by the unique index; a rewritten core is refused by the append-only triggers; a coercer who is the institution is not stopped, and the record says so.}
\qa{Checked by}{\code{check\_c1c10\_objects\_resolve}, the C3 property tests and the C3 model, the concurrency suite, \code{check\_card\_profile}, \code{check\_card\_emulator}, \code{check\_verifier\_device}, \code{check\_duress\_on\_card}, \code{check\_duress\_claims\_are\_aware}, the attack suite's forgery and tamper cases.}
\qa{Now or later}{Now. It is the part of Polaris that has existed since Version 1; the card is an object now and silicon is later.}
\qa{Inside and outside}{Inside: the person, who exists outside the software. Outside: the wall of constraints (Section~\ref{sec:walls}).}
\end{layercard}

\analogy{The town hall keeps one register. Each entry says who a person is, which official entered
it, under which seal, and what papers the official was shown, though never the papers themselves. The
register cannot be edited, only appended to; a second entry for the same person is refused at the
desk; and the seal is one that every reader can check against the official's published mark without
asking the hall.}
